\scanpage{15}
\mspage{\textcircled{2}}
Soit maintenant $k$ une ext. \uuline{alg. close} de $\mathbb Q$ (pas néc. alg/$\mathbb Q$), et $X$ une courbe de type $g,\nu$ sur $k$, $x\in X$. Alors $X$ correspond à un pt géom
\[
\xi\in\operatorname{Ob}\uline{\mathrm{Pt}}(M_{g,\nu},k)\simeq\operatorname{Ob}\uline{\mathrm{Pt}}(M_{g,\nu k}/k)
\]
\[
X\simeq(\mathcal X_{g,\nu k})_\xi
\]
\[
\deleted{\unsure{\xi}},\ x\in\operatorname{Ob}\uline{\mathrm{Pt}}(\mathcal X_{g,\nu\unsure{k}},k)\simeq\operatorname{Ob}\uline{\mathrm{Pt}}(M_{g,\nu k}/k)\qquad(x\text{ au dessus de }\xi)
\]

\deleted{$\unsure{\Pi}$} Considérons
\deleted{$\unsure{\Pi=\pi_1(X,x)}$}
\[
(11)\quad\left\{\begin{array}{ll}
\Gamma_X\overset{\mathrm{déf}}=\pi_1(M_{g,\nu k},\xi),&\mathcal M_X=\pi_1(M_{g,\nu},\xi)\\
\mathcal G_{X,x}=\pi_1(\mathcal X_{g,\nu k},x),&\mathcal N_{X,x}=\pi_1(\mathcal X_{g,\nu},x)\\
\Pi_{X,x}=\pi_1(X,x)&
\end{array}\right.
\]
d'où
\diagram{\begin{tikzcd}[column sep=small]
&&&1\arrow[d]&1\arrow[d]&\\
(12)\quad&1\arrow[r,no head]&\Pi_{X,x}\arrow[r]\arrow[d,equal]&\mathcal G_{X,x}\arrow[r]\arrow[d]&\Gamma_X\arrow[r]\arrow[d,no head]&1\\
&1\arrow[r]&\Pi_{X,x}\arrow[r]&\mathcal M_{X,x}\arrow[r,no head]\arrow[d]&\mathcal N_X\arrow[r,no head]\arrow[d]&1\\
&&&\Gamma_{\widetilde{\mathbb Q}/\mathbb Q}\arrow[r,"\simeq"]&\Gamma_{\widetilde{\mathbb Q}/\mathbb Q}\arrow[d]&\deleted{\unsure{1}}\\
&&&&1&
\end{tikzcd}}
\margin{où $\widetilde{\mathbb Q}$ est la clôture alg. de $\mathbb Q$ dans $k$}
\deleted{\unsure{On a des iso extérieurs}}

\deleted{\unsure{On}} On a un isom. extérieur can. \deleted{$\unsure{\mathcal G_{X,x}\simeq\widehat{\mathcal G}_{g,\nu}}$} \deleted{\unsure{les isom sont définis}}
\[
(13)\qquad\mathcal M_{X,x}\xrightarrow{\sim}\mathcal M_{g,\nu}
\]
i.e. isom, définis à autom. int. près, définissant un iso de suites exactes entre (12) et (5).

Conséquence :
\diagram{\deleted{\begin{tikzcd}
(14)\quad1\arrow[r,no head]&\Pi_{X,x}\arrow[r,no head]\arrow[d,"\simeq"']&\mathcal G_{X,x}\arrow[r,no head]\arrow[d,"\simeq"']&\Gamma_X\arrow[r]\arrow[d,"\simeq"']&1\\
1\arrow[r,no head]&\widehat\Pi_{g,\nu}\arrow[r,no head]&\widehat{\mathcal G}_{g,\nu}\arrow[r,no head]&\widehat\Gamma_{g,\nu}\arrow[r,no head]&1
\end{tikzcd}}}
\[
(14)\qquad\Pi_{X,x}\simeq\widehat\Pi_{g,\nu}
\]
\inserted{iso} défini \uline{mod opération de}
\[
\mathcal M_{g,\nu}\supset\widehat{\mathcal G}_{g,\nu}
\]
C'est une structure supplémentaire \inserted{can.} sur le groupe profini extérieur $\Pi_{X,x}$, appelé \deleted{\unsure{son}} ``arithmétisation canonique'' : restriction du groupe structural

\mspage{14}
